\documentclass[aspectratio=169,10pt,spanish]{beamer}

\input{../../../_preambulo_beamer.tex}
\input{../../../_comandos_beamer.tex}

\selectlanguage{spanish}

\title{MA1001B - Modelación Estadística para la Toma de Decisiones}
\subtitle{Lección 26: Estimación puntual: insesgadez, eficiencia, consistencia}
\author{}
\institute{Tecnol\'ogico de Monterrey}
\date{}

\begin{document}

\begin{frame}[plain]
  \vspace{1.2cm}
  {\Large\bfseries MA1001B - Modelación Estadística para la Toma de Decisiones\par}
  \vspace{0.7cm}
  {\large Lección 26: Propiedades de la estimación puntual\par}
  \vspace{0.7cm}
  {\color{TecRojo}\rule{\textwidth}{1.2pt}}
  \vspace{0.8cm}

  Facultad MA1001B\\[0.3cm]
  Periodo\\[0.3cm]
  Tecnol\'ogico de Monterrey
\end{frame}

\begin{frame}{La pregunta del estimador}
  \begin{alertblock}{Pregunta guía}
    Si los costos son $N(100,15)$, ¿por qué preferir $\bar x$ a una sola
    observación $X_1$ como estimador puntual de $\mu$?
  \end{alertblock}

  \vspace{0.6em}
  Al final de esta lección, usted puede:
  \begin{itemize}
    \item definir la insesgadez como $E[\hat\theta]=\theta$;
    \item comparar la eficiencia mediante las varianzas de los estimadores;
    \item reconocer la consistencia como la reducción de la dispersión con $n$;
    \item explicar por qué la insesgadez no implica que una muestra pequeña sea precisa.
  \end{itemize}

  \vfill
  \scriptsize Todos los costos usados hoy son sintéticos. No se usan datos reales de empresa.
\end{frame}

\begin{frame}{Dos estimadores insesgados}
  \small
  Población $X\sim N(100,15^2)$. Ambos
  \[
    E[\bar x]=\mu\qquad\text{y}\qquad E[X_1]=\mu.
  \]
  8{,}000 muestras con semilla 42 de $n=20$:
  \[
    \text{media de }\bar x=99.957310,\quad \text{sesgo}=-0.042690,
  \]
  \[
    \text{media de }X_1=99.916074,\quad \text{sesgo}=-0.083926.
  \]

  \vspace{0.3em}
  \begin{exampleblock}{Insesgadez}
    Correcto en promedio no es lo mismo que varianza pequeña.
  \end{exampleblock}
\end{frame}

\begin{frame}[fragile]{Paso 1: Insesgadez}
  \begin{columns}[T]
    \begin{column}{0.42\textwidth}
      \small
      \begin{lstlisting}
xbar = draws.mean(axis=1)
one = draws[:, 0]
bias_xbar = xbar.mean() - 100
      \end{lstlisting}

      \begin{block}{Salida verificada}
        Media de $\bar x=99.957310$\\
        Media de $X_1=99.916074$
      \end{block}
    \end{column}
    \begin{column}{0.58\textwidth}
      \centering
      \includegraphics[width=\textwidth]{../figuras/point_estimation_01_unbiasedness.png}
      \vspace{0.2em}
      \scriptsize Medias muestrales del Paso 1
    \end{column}
  \end{columns}
\end{frame}

\begin{frame}{Eficiencia y consistencia}
  \small
  En $n=20$,
  \[
    \mathrm{Var}(X_1)=225,\qquad \mathrm{Var}(\bar x)=11.25.
  \]
  Valores empíricos: $227.612313$ versus $11.242349$, razón $20.245975$.
  Subir $n$ a 200 reduce $\mathrm{SD}(\bar x)$ a $1.052019$ versus la
  fórmula $15/\sqrt{200}=1.060660$.
\end{frame}

\begin{frame}[fragile]{Paso 2: Comparación de dispersión}
  \begin{columns}[T]
    \begin{column}{0.40\textwidth}
      \small
      \begin{lstlisting}
var_xbar = xbar.var()
var_one = one.var()
ratio = var_one / var_xbar
      \end{lstlisting}

      \begin{block}{Salida verificada}
        $\mathrm{SD}(\bar x_{20})=3.352961$\\
        $\mathrm{SD}(\bar x_{200})=1.052019$
      \end{block}
    \end{column}
    \begin{column}{0.60\textwidth}
      \centering
      \includegraphics[width=\textwidth]{../figuras/point_estimation_02_efficiency_consistency.png}
      \vspace{0.2em}
      \scriptsize Eficiencia y consistencia del Paso 2
    \end{column}
  \end{columns}
\end{frame}

\begin{frame}{Paso 3: Insesgado pero demasiado ruidoso}
  \begin{columns}[T]
    \begin{column}{0.42\textwidth}
      \small
      En $n=5$, $\bar x$ sigue siendo insesgado:
      media $100.066209$, SD $6.737762$.

      \vspace{0.4em}
      $P(|\bar x-100|>10)=0.137875$.\\
      En $n=200$ ese evento es $0.000000$.

      \vspace{0.4em}
      \textbf{Límite:} la insesgadez no implica varianza pequeña.
      La Lección 27 envuelve $\bar x$ en un intervalo $z$.
    \end{column}
    \begin{column}{0.58\textwidth}
      \centering
      \includegraphics[width=\textwidth]{../figuras/point_estimation_03_small_sample_variance.png}
      \vspace{0.2em}
      \scriptsize Ruido de muestra pequeña del Paso 3
    \end{column}
  \end{columns}
\end{frame}

\begin{frame}{Verificación de conceptos}
  \begin{enumerate}
    \item ¿Por qué $\bar x$ y $X_1$ son ambos insesgados para $\mu$?
    \item ¿Qué dice la razón $\mathrm{Var}(X_1)/\mathrm{Var}(\bar x)\approx 20$
      en $n=20$?
    \item ¿Cómo cambia $\mathrm{SD}(\bar x)$ de $n=20$ a $n=200$?
    \item ¿Por qué una media insesgada con $n=5$ puede seguir siendo la
      herramienta de decisión equivocada?
  \end{enumerate}

  \vspace{0.6em}
  \begin{block}{Conexión con la Lección 27}
    La sesión puede continuar directamente con un intervalo de confianza $z$
    para una media cuando $\sigma$ es conocida.
  \end{block}

  \vfill
  \scriptsize Fuente: OpenStax, \textit{Introductory Business Statistics 2e},
  introducción del Capítulo 8. CC BY-NC-SA.
\end{frame}

\appendix



\end{document}
